\chapter*{Preface}
\addcontentsline{toc}{chapter}{Preface}

\fml{} is a C++17 library for cosmology and the large-scale structure of the Universe. It started as a collection of
the tools needed for fast approximate simulations (particle--mesh gravity, Lagrangian perturbation theory, the COLA
method and modified-gravity solvers) and has grown into a library that covers the whole chain from the expansion
history and the linear Einstein--Boltzmann equations, through initial conditions and $N$-body, hydrodynamical and
wave dark matter simulations, to halo and void finding, clustering statistics, survey estimators, analytic nonlinear
models and statistical inference.

This book documents all of it. For every feature it gives
\begin{itemize}
\item the equations that are solved, derived to the extent needed to state the conventions unambiguously;
\item the algorithms used to solve them, with the numerical choices that matter for accuracy and cost;
\item how the feature is used from code (the main classes and functions, a short example and the important options);
\item how it is validated, and against which external codes and papers;
\item references to the literature and to the codes the methods come from or are compared with.
\end{itemize}
The module guides in the repository (\file{README.md} files next to the code) remain the authoritative description of
interfaces and test procedures; this book collects the physics and numerics behind them in one place, with a common
notation. \Cref{app:module-index} maps every module directory of the repository to the chapters that describe it.

\paragraph*{How to read this book.} \Cref{part:intro} (\cref{ch:introduction,ch:conventions,ch:getting-started}) gives an overview, shows how to build and run the code and
fixes the conventions (units, Fourier transforms, grids, parallel layout, cosmological notation) used everywhere
else. The remaining parts follow the flow of a typical calculation: the numerical infrastructure, the homogeneous
universe and linear theory, analytic models of nonlinear structure, initial conditions, simulations, structure
identification, clustering statistics, and finally surveys, lensing and inference. Each chapter starts with a summary
of the modules it covers and can be read on its own.

\paragraph*{Status.} The book is written alongside the code and describes the library as of the date under each
chapter title. Facts, defaults and validation numbers are taken from the source, the module documentation and the
tests; where a result is a one-off measurement or an experiment rather than a regression test, the text says so. When
the code changes, the module's own documentation is the authority until the chapter is updated.
